\documentclass[11pt]{article}
\usepackage[a4paper,margin=2.2cm]{geometry}
\usepackage[T1]{fontenc}
\usepackage{lmodern}
\usepackage{microtype}
\usepackage{amsmath}
\usepackage{graphicx}
\usepackage{booktabs}
\usepackage{adjustbox}
\usepackage{longtable}
\usepackage{float}
\usepackage{caption}
\usepackage{listings}
\usepackage[numbers,sort&compress]{natbib}
\usepackage[hidelinks]{hyperref}
\lstset{basicstyle=\scriptsize\ttfamily,breaklines=true,columns=fullflexible,keepspaces=true}
\captionsetup{font=small}
\input{numbers.tex}
% A figure is skipped when the analysis did not make it (e.g. Phase 2 figures when the gate stopped the study).
\newcommand{\fig}[3]{\IfFileExists{figures/#1.pdf}{\begin{figure}[htbp]\centering\includegraphics[width=\textwidth]{figures/#1.pdf}\caption{#2}\label{#3}\end{figure}}{}}

\title{Right for the Right Reasons? Firearm Grounding Feedback in VERA's Verbalized Learning}
\author{VISILAB firearm-action dataset, frozen InternVL2-8B}
\date{\today}

\begin{document}
\maketitle

\begin{abstract}
VERA adapts a frozen vision-language model by rewriting the questions it answers, and scores only the answer text. We asked whether InternVL2-8B, adapted this way to the VISILAB firearm-action dataset, decides a clip's category from the firearm, and whether telling VERA's optimizer where the model looked would make it rely on the firearm more. Two pre-registered measures were taken on all 398 videos with the untrained questions: the decision's attention to the annotated firearm relative to chance (lift), and the drop in the true category's probability when the firearm is greyed out rather than an equal patch elsewhere (reliance). The attention is specific to the firearm: lift exceeds that of another video's boxes by \val{gate_V1_estimate} [\val{gate_V1_low}, \val{gate_V1_high}]. Correct answers depend on the firearm entirely: every one of them changes when it is masked. But lift does not track reliance ($\rho=$ \val{gate_V2_estimate} [\val{gate_V2_low}, \val{gate_V2_high}]), so the pre-registered gate failed and the training experiment did not run. The model looks at handguns yet names only \val{R5_seed_Handgun_category_acc_pct}\,\% of the Handgun videos correctly, calling most of the rest No\_Gun: its failure is recognition, not where it looks, and attention feedback would not address it.

\end{abstract}

\section{Questions}
Training~1 adapted VERA~\cite{ye2025vera} to the VISILAB firearm-action dataset~\cite{ruizsantaquiteria2024firearm}: a frozen InternVL2-8B~\cite{chen2024internvl15} rewrites its own guiding questions from the errors it makes, and only the answer text is scored. An answer can be right for the wrong reasons~\cite{ross2017right}: the category Handgun can come from the room, the actor or the pose rather than from the gun. The dataset's firearm boxes let us look beyond the text. This study asks two questions.
\begin{enumerate}
\item \textbf{RQ1 (measurement).} When the frozen model decides a clip's category, does it look at the annotated firearm (attention), does it rely on it (occlusion), and do the two agree?
\item \textbf{RQ2 (training).} If VERA's optimizer is also told where the model looked, do the learned questions make the model rely more on the firearm, without losing more than 5 points of category accuracy?
\end{enumerate}
The design, the gate and the outcomes were fixed in a pre-registration before any grounding number was taken; the one deviation (D1) is in Section~\ref{sec:deviations}.

\section{Setting}
Everything that is not about grounding is Training~1 unchanged: the 240/62/96 split by recording group, the codebook and learner prompt, eight frames per video sampled evenly over the whole video, InternVL2-8B revision 6fb9ad69 in BF16 with greedy decoding and 1,024 new tokens, batches of two videos, ten epochs and validation every 100 batches within an epoch. Boxes come from the dataset's annotations after the audit of the earlier study. A sampled frame is \emph{boxed} when its audited status is positive, and a video is \emph{grounding-evaluable} when at least one of its eight frames is boxed; No\_Gun videos never are. Training reads training-split boxes only, selection reads validation boxes, and test boxes are read only for the final evaluation.

\section{Measures}
All measures are taken on the model's own greedy answer under the questions being evaluated.

\paragraph{Decision.} The answer is cut at its last \texttt{Category:} field. The three category names are scored as continuations of that prefix (summed token log-probabilities, normalised over the three, FP32 head), giving $p_\text{true}$. The decision position is the last token before the category name.

\paragraph{M1, attention lift (correlational).} One scoring forward pass with eager attention keeps the decision position's attention row in every layer. Frame $k$'s image tokens are the $k$-th run of 256 tokens in the single image span; token $(r,c)$ covers a $62.5\times62.5$ cell of the 0--1000 grid of the $448\times448$ input. With $a$ the frame's attention (mean over 32 layers and 32 heads, renormalised within the frame) and $c$ each token's fractional box coverage, the video lift is $\sum_k a_k\!\cdot\!c_k \,/\, \sum_k \lVert c_k\rVert_1/256$ over boxed frames: 1.0 is chance. Sensitivities: layer groups 0--10, 11--21 and 22--31, and removing each frame's single most-attended token before renormalising (attention sinks~\cite{xiao2024sinks,darcet2024registers}).

\paragraph{M2, occlusion reliance (causal).} Following occlusion analysis~\cite{zeiler2014visualizing}, the answer is regenerated twice: once with every boxed frame's firearm box filled with grey (127), once with an equal-sized grey box at a control position (horizontal mirror, else vertical mirror, else the disjoint position closest to the centre). Reliance $=\log p_\text{true}(\text{control}) - \log p_\text{true}(\text{box})$, in nats: positive when hiding the gun hurts the true answer more than hiding an equal patch elsewhere. Frames whose box has no disjoint control position are left out of both masks.

\paragraph{M3 and M4.} Category, action-group and code accuracy with Training~1's strict parser, and the share of answers whose explanation names a weapon (on No\_Gun videos, a false mention). A grounding-evaluable video with the correct category and reliance $\le 0$ is \emph{right for the wrong reasons}.

\section{Design}
\paragraph{Phase 0.} Real-model checks before any grounding number: our generation equals the runner's \texttt{batch\_chat}; an eager replay reproduces the greedy tokens; eager and FlashAttention category scores agree within 0.05 nats; greying one cell changes the vision token at that cell (token-grid orientation).

\paragraph{Phase 1 and the gate.} The seed questions on all 398 videos. V1 (spatial specificity): lift with the true boxes minus lift with \emph{shuffled} boxes (another grounding-evaluable video of the same camera, frame by frame). V2 (agreement): Spearman $\rho$ between lift and reliance. Phase~2 runs only if both 95\,\% intervals lie above 0; attention is otherwise not a valid training signal and the study stops at RQ1, because attention and reliance need not agree~\cite{jain2019attention}.

\paragraph{Phase 2.} Three arms with three training-order seeds each (20261026, 20261027, 20261028), paired by seed. \textbf{T1} is Training~1 (its own run is seed 20261026). \textbf{G} adds a block to the optimizer prompt after the targets: for each video in the batch, the attention lift on the annotated firearm in each boxed frame, measured on the learner's own answer, with an overall level (low/typical/high at Phase~1's training-split quartiles) and one fixed sentence asking for questions that lead the model to look at the firearm without lowering accuracy (Appendix~\ref{app:prompts}). \textbf{P} (placebo) shows the identical block with each frame's value drawn at random from Phase~1's training-split frame lifts, so that any effect of G beyond P comes from the measured values, not from the block's presence or its sentence. T1 selects by validation category accuracy; G and P select, among validated steps within 5 points of their best validation category accuracy, the step with the highest pooled validation lift. Training feedback is attention (M1) while the primary outcome is occlusion (M2), so questions that move attention without changing reliance cannot pass.

\paragraph{Outcomes.} On the test split: \textbf{H1 (primary)} mean reliance of G exceeds P, as the per-video G$-$P difference averaged over seeds; \textbf{H2 (co-primary)} category accuracy of G is at least T1's minus 5 points (lower bound of the paired difference above $-0.05$). Intervals are 95\,\% percentile intervals from 2,000 bootstrap draws resampling test recording groups within category (seed 20261026)~\cite{efron1993bootstrap}. Secondary contrasts are descriptive and not corrected for multiplicity.

\section{Results}
% PROSE PENDING: written from the saved analysis after the run. Floats below are final.
\subsection{Measurement checks}
\input{tables/R1.tex}

\subsection{RQ1: does the frozen model look at and rely on the firearm?}
\input{tables/R2.tex}
\fig{F1_phase1_lift_and_reliance}{Seed questions, all splits. Left: attention lift with the true and with shuffled boxes (log scale, 1 = chance). Right: lift against occlusion reliance per grounding-evaluable video.}{fig:F1}
\input{tables/R2s.tex}

\subsection{RQ2: grounding feedback in training}
\input{tables/R3.tex}
\fig{F3_validation_trajectories}{Validation category accuracy (left) and pooled validation lift (right, G and P) at each validation step; one line per seed.}{fig:F3}
\input{tables/R4.tex}
\fig{F2_arms}{Test reliance, lift and category accuracy per arm. Dots: seeds; bars: mean over seeds; dashed: seed questions.}{fig:F2}
\input{tables/R5.tex}
\input{tables/R6.tex}
\input{tables/R7.tex}


\section{Deviations}\label{sec:deviations}
\textbf{D1} (after Phase~0, before Phase~1). Phase~0's checks b2 (eager replay reproduces 207 of 209 greedy tokens) and b3 (largest eager--FlashAttention gap 0.12 nats) failed their strict thresholds. A diagnostic found the same two token mismatches in a FlashAttention replay, at near-ties (margins $-0.19$ and $-0.42$), and the 0.12-nat gap on a class at $p\approx0.002$; for the chosen class the gap was at most 0.0007 nats. Attention is therefore read from the eager pass and every probability from a FlashAttention pass, the kernels that generate the answers. Nothing else changed.

\section{Limitations}
\begin{itemize}
\item \textbf{One frozen model and one question set.} RQ1 is answered for InternVL2-8B with the seed questions. Learned questions could change where the model looks and how it decides; that was Phase~2's question.
\item \textbf{Lift depends on box size.} The ceiling of $256/\text{covered tokens}$ makes lifts of small and large firearms hard to compare, and it lowers the pooled correlation with reliance. The pre-registered gate pools sizes; the size-stratified correlations are descriptive and were not pre-registered.
\item \textbf{Reliance is nearly binary.} With the seed questions, a correct firearm answer flips whenever the firearm is masked, and a wrong one rarely changes. Reliance therefore measures mostly whether a firearm was recognised, and its correlation with any other measure inherits that.
\item \textbf{Masks and boxes.} Grey patches create edges, which the equal-sized control handles; but \val{phase1_controls_within_one_cell_pct}\,\% of control patches lie within one token cell of a firearm and may cover part of the hand or firearm, pulling reliance toward 0. Boxes are interpolated between keyframes and can be loose; only audited positive frames count.
\item \textbf{Eight frames in one image block.} Each frame's 256 tokens are located exactly, but the model can mix information across frames, so per-frame attention is not per-frame use.
\item \textbf{Weapon mentions} (M4) count denials such as ``not holding a firearm'' as mentions, as pre-registered.
\item \textbf{The dataset} has four actors, mostly one viewpoint per room and scripted actions; a model that misses handguns here may behave differently on surveillance footage.
\item \textbf{Disclosed look at partial data.} The gate was computed on the first 72 videos while testing this report, before the review-driven fixes and the restart (Section~\ref{sec:deviations}).
\end{itemize}


\bibliographystyle{unsrtnat}
\bibliography{references}

\appendix
\section{Prompts}\label{app:prompts}
\subsection*{Optimizer instruction of G and P}
Training~1's optimizer instruction with the grounding block inserted after the targets; T1 uses it without the block.
\lstinputlisting{optimizer_grounded.txt}
\subsection*{A grounding block from G's training}
\lstinputlisting{grounding_example.txt}

\section{Questions}\label{app:questions}
\input{questions.tex}
\end{document}
